\chapter{Marchés monétaires}\label{ch:m2:money-markets}

Le jeudi 29 décembre 2022, les \omterm{def:m2:money-markets:mmf}{fonds monétaires} et les autres
établissements éligibles ont déposé pour une nuit 2\,308~milliards de dollars auprès de la Réserve fédérale
via sa \omterm{def:m2:money-markets:rrp}{facilité de reverse repo}. Le vendredi 30 décembre, le dernier jour ouvré
de l'année, ils ont déposé 2\,554~milliards de dollars, soit 245~milliards de plus
en un jour, un record. Le premier jour ouvré de 2023, le montant est retombé à 2\,188
milliards de dollars. Rien dans la politique monétaire n'a changé au cours de ces trois jours. Ce qui
a changé, c'est le calendrier~: le dernier jour de l'année, les banques et les dealers
qui empruntent d'ordinaire ces liquidités refusaient de les prendre, car la taille de leur
bilan cette nuit-là est ce que leurs régulateurs et leurs investisseurs
mesurent. Les marchés monétaires, où les liquidités sont prêtées pour un jour à un an, comptent
parmi les marchés les plus profonds qui soient, et ils sont façonnés autant par les dates
comptables que par les taux d'intérêt. Ce chapitre traite de leurs instruments, de leurs
conventions de cotation et des dates auxquelles ils dysfonctionnent.

\section{Bons du Trésor et leurs conventions de cotation}

\begin{definition}[Bon du Trésor]\label{def:m2:money-markets:bill}
Un \emph{bon du Trésor}\index{bon du Trésor} (Treasury bill) est un titre de dette publique
à échéance d'un an au plus, qui ne verse aucun coupon~: il est vendu en dessous de sa
valeur nominale et rembourse cette valeur à l'échéance. Le Trésor américain émet par adjudication des bons
de 4, 6, 8, 13, 17, 26 et 52 semaines, ainsi que des bons de gestion de trésorerie de durée
irrégulière, lorsque ses besoins de trésorerie l'exigent.
\end{definition}

Le prix d'un bon s'exprime par un taux, et le marché en utilise trois, qui
diffèrent par la base sur laquelle l'intérêt est exprimé et par la durée de l'année.

\begin{definition}[Taux d'escompte, taux monétaire et rendement équivalent obligataire]\label{def:m2:money-markets:yields}
Pour un bon ayant $t$ jours à courir et un prix $P$ pour 100 de valeur nominale~:
le \emph{taux d'escompte}\index{taux d'escompte} (discount yield) $d$, calculé en base d'escompte bancaire, est
la décote rapportée à la valeur nominale sur une année de 360~jours,
$P = 100\,(1 - d\,t/360)$~; le \emph{taux monétaire}\index{taux monétaire}
(money-market yield) $m = \frac{100 - P}{P}\frac{360}{t}$ est l'intérêt simple sur le montant
versé sur une année de 360~jours, soit le taux d'un dépôt de même durée~; le
\emph{rendement équivalent obligataire}\index{rendement équivalent obligataire} (bond-equivalent yield) $i$ (le
\emph{taux de placement} du Trésor) est le même sur une année de 365~jours, $i = \frac{100 -
P}{P}\frac{y}{t}$ avec $y = 365$, ou 366 si l'année suivant l'émission contient un
29~février, pour les bons de six mois au plus~; au-delà, elle vérifie
\[
P\Bigl[1 + \Bigl(t - \frac y2\Bigr)\frac{i}{y}\Bigr]\Bigl(1 + \frac i2\Bigr) = 100,
\]
une capitalisation semestrielle qui la rend comparable au rendement d'une
obligation à coupon (voir du \cref{ch:m2:government-bonds}).
\end{definition}

\begin{proposition}[Ordre des trois rendements]\label{prop:m2:money-markets:order}
Pour $0 < t \le y/2$ et $d > 0$,
\[
m = \frac{d}{1 - d\,t/360} > d, \qquad i = m\,\frac{y}{360} > m .
\]
Le \omterm{def:m2:money-markets:yields}{taux d'escompte} sous-estime toujours le rendement des liquidités investies~; l'écart
croît avec le taux et avec la durée.
\end{proposition}
\begin{proof}
De $P = 100(1 - dt/360)$, on tire $100 - P = 100\,dt/360$, donc $m = \frac{dt/360}{1 -
dt/360}\frac{360}{t} = \frac{d}{1 - dt/360}$, supérieur à $d$ car le
dénominateur est inférieur à un. La seconde égalité est la définition avec $y$ à la
place de 360.
\end{proof}

\begin{example}[Un bon à treize semaines]\label{ex:m2:money-markets:13w}
Un bon de 91 jours adjugé à un \omterm{def:m2:money-markets:yields}{taux d'escompte} de 3.82~\% coûte
$100(1 - 0.0382 \times 91/360) = 99.034389$ pour 100. Son \omterm{def:m2:money-markets:yields}{taux monétaire} est
3.8572~\% et son \omterm{def:m2:money-markets:yields}{rendement équivalent obligataire} 3.9108~\%~: neuf points de base séparent
le chiffre affiché à l'écran d'adjudication et celui à comparer avec un
obligation du Trésor. Les formules sont celles du Trésor lui-même, et elles reproduisent ses
exemples publiés au dernier chiffre près (voir la \cref{bld:m2:money-markets:mmyield}).
\end{example}

\begin{omfigure}
\begin{tikzpicture}
\begin{axis}[width=11cm,height=5cm,
  xlabel={échéance (semaines)},ylabel={taux (\%)},
  tick label style={font=\small},label style={font=\small},
  xmin=0,xmax=54,ymin=3.68,ymax=3.94,grid=major,grid style={gray!25},
  xtick={4,8,13,17,26,52},
  legend columns=3,legend style={font=\footnotesize,at={(0.5,-0.3)},anchor=north,draw=none,fill=none,/tikz/every even column/.append style={column sep=8pt}},
  table/col sep=comma]
\addplot[omDef,very thick,mark=*,mark size=1.5pt] table[x=weeks,y=discount]{figdata/markets-2/02-money-markets/bills.csv};
\addplot[omProp,very thick,mark=square*,mark size=1.5pt] table[x=weeks,y=mmy]{figdata/markets-2/02-money-markets/bills.csv};
\addplot[omThm,very thick,mark=triangle*,mark size=2pt] table[x=weeks,y=bey]{figdata/markets-2/02-money-markets/bills.csv};
\legend{taux d'escompte,taux monétaire,rendement équivalent obligataire}
\end{axis}
\end{tikzpicture}

\omcaption{Une courbe de bons illustrative (les sept échéances adjugées) selon
les trois conventions. Les \omterm{def:m2:money-markets:yields}{taux d'escompte} diminuent régulièrement avec l'échéance~; les \omterm{def:m2:money-markets:yields}{taux
monétaires} ne le font pas, car l'écart entre eux s'élargit avec la durée~: le
bon à 52 semaines rapporte \emph{davantage} que celui à 26 semaines dans cette convention. Seule la
courbe du \omterm{def:m2:money-markets:yields}{rendement équivalent obligataire} peut être raccordée à la courbe des coupons. Les niveaux sont
indicatifs, et non des cotations. Données~: celles du
tutoriel du chapitre.}\label{fig:m2:money-markets:bills}
\end{omfigure}

\section{Billets de trésorerie et certificats de dépôt}

\begin{definition}[Billet de trésorerie]\label{def:m2:money-markets:cp}
Un \emph{billet de trésorerie}\index{billet de trésorerie} (commercial paper) est une dette à court terme
non garantie émise par des sociétés, des banques et des entités ad hoc (SPV) auprès
d'investisseurs, en général à escompte. Aux États-Unis, il est dispensé d'enregistrement si son
échéance n'excède pas 270 jours et qu'il finance des opérations courantes~; la plupart
est émise pour des durées bien plus courtes et renouvelée en permanence.
Le \emph{papier adossé à des actifs} est émis par une entité ad hoc qui détient des créances et
rembourse à partir de celles-ci.
\end{definition}

\begin{definition}[Certificat de dépôt]\label{def:m2:money-markets:cd}
Un \emph{certificat de dépôt}\index{certificat de dépôt} (certificate of deposit, CD) est un dépôt bancaire
à terme fixe, émis comme titre négociable~: la banque verse
le principal et un intérêt simple à l'échéance, et le porteur peut revendre le
certificat avant cette échéance.
\end{definition}

Les \omterm{def:m2:money-markets:cp}{billets de trésorerie} et les \omterm{def:m2:money-markets:cd}{certificats de dépôt} rapportent un spread au-dessus des \omterm{def:m2:money-markets:bill}{bons du Trésor} pour trois raisons~: le crédit (une entreprise ou
une banque peut faire défaut, alors qu'un gouvernement qui emprunte dans sa propre monnaie ne peut pas
faire défaut de la même façon), la liquidité (il n'existe pas de marché secondaire profond pour le papier
à trente jours d'un même émetteur), et la base d'investisseurs (un \omterm{def:m2:money-markets:bill}{bon du Trésor} peut être détenu par n'importe quel
\omterm{def:m2:money-markets:mmf}{fonds monétaire}, le papier seulement par certains). Le spread est faible sur les marchés calmes et c'est le
premier élément à bouger quand le calme prend fin, car un porteur qui ne peut pas vendre le papier
doit attendre son échéance, et un émetteur qui ne peut pas renouveler son papier doit trouver
les liquidités ailleurs, auprès d'une ligne de crédit bancaire qu'il a négociée à cet effet.

\begin{method}[Comparer les instruments courts]\label{met:m2:money-markets:compare}
Convertir chaque cotation en \omterm{def:m2:money-markets:yields}{taux monétaire} (intérêt simple sur le montant
payé, actual/360) avant toute comparaison~: le \omterm{def:m2:money-markets:yields}{taux d'escompte} d'un bon, le \omterm{def:m2:money-markets:yields}{taux d'escompte}
d'un papier, le coupon d'un \omterm{def:m2:money-markets:cd}{certificat de dépôt}, un taux de repo. Comparer alors des échéances égales, et se demander
pour chaque spread lequel du crédit, de la liquidité ou de la base d'investisseurs il rémunère.
\end{method}

\section{Fonds monétaires}

\begin{definition}[Fonds monétaire]\label{def:m2:money-markets:mmf}
Un \emph{fonds monétaire}\index{fonds monétaire} (money-market fund, MMF) est un fonds commun de placement
restreint à des instruments courts et de haute qualité, offrant une liquidité quotidienne et
une valeur de la part qui varie peu, voire pas du tout. Les fonds \emph{gouvernementaux} détiennent
des \omterm{def:m2:money-markets:bill}{bons du Trésor}, d'autres titres publics et des repos adossés à ceux-ci~; les fonds \emph{prime}
détiennent aussi du papier et des dépôts bancaires~; beaucoup de fonds gouvernementaux et de fonds de détail
maintiennent une valeur de part stable à un dollar.
\end{definition}

Un fonds qui promet un dollar stable et investit dans tout ce qui peut perdre
de la valeur fait une promesse qu'il ne peut pas tenir. Le 16 septembre 2008, le lendemain
du dépôt de bilan de Lehman Brothers, le Reserve Primary Fund, un fonds prime de 62~milliards de dollars
qui détenait 785~millions de dollars de dette Lehman et a valorisé cette dette à zéro
et sa part à 97~cents. Les investisseurs ont retiré environ 300~milliards de dollars, soit 14~\%,
des fonds prime cette semaine-là~; les fonds ont vendu du papier pour les rembourser, le marché du papier
s'est arrêté, et le Trésor américain a garanti la valeur de la part en dollars de plus de 3
billions de dollars de parts de fonds. Chaque réforme des \omterm{def:m2:money-markets:mmf}{fonds monétaires} depuis lors, la plus récente
adoptée par le régulateur américain en 2023, a porté sur les fonds autorisés à promettre un
prix stable et sur ce qu'ils doivent faire quand les investisseurs sortent plus vite que
leurs actifs n'arrivent à échéance.

\begin{dated}{2026-09}{Fonds monétaires américains}\label{dat:m2:money-markets:mmf}
\textbf{Taille} (semaine au 16 septembre 2026)~: 7.92~billions de dollars, dont
les fonds gouvernementaux 6.53 billions, les fonds prime 1.24 billions et les fonds exonérés d'impôt
0.15~billion~; fonds de détail 3.11 billions, institutionnels 4.81 billions.
\textbf{Règles} (réformes adoptées en juillet 2023, en vigueur d'ici octobre 2024)~: au moins
25~\% des actifs liquides en un jour et 50~\% en une semaine (auparavant 10~\% et
30~\%)~; plafonnements des rachats supprimés~; les fonds prime institutionnels et les
fonds institutionnels exonérés d'impôt doivent prélever des frais de liquidité les jours où les rachats nets
dépassent 5~\% des actifs nets~; les autres fonds non gouvernementaux peuvent prélever jusqu'à 2~\%.
\end{dated}

\begin{definition}[Facilité de reverse repo]\label{def:m2:money-markets:rrp}
La \emph{facilité de reverse repo}\index{facilité de reverse repo} d'une banque centrale permet
aux contreparties éligibles de lui prêter des liquidités au jour le jour contre des titres, à un taux
qu'elle fixe~; ce taux place un plancher sous ce que ces contreparties
accepteront de quiconque d'autre.
\end{definition}

La \cref{fig:m2:money-markets:plumbing} montre où va l'argent. Un fonds gouvernemental
place ses liquidités en trois endroits~: \omterm{def:m2:money-markets:bill}{bons du Trésor}, repo avec les dealers et, lorsque
les dealers refusent de la prendre, la \omterm{def:m2:money-markets:rrp}{facilité de reverse repo} de la Federal Reserve
(voir la \cref{rem:m2:central-banks-and-the-short-rate:below}). Lequel des trois
l'emporte chaque jour est décidé par leurs taux au point de base près, et le choix
du fonds fait bouger le coût de financement du Trésor et celui des dealers.

\begin{omfigure}
\begin{tikzpicture}[font=\small,
  box/.style={draw,thick,rounded corners=2pt,align=center,minimum height=1.0cm,inner sep=4pt},
  arr/.style={-{Stealth[length=2.5mm]},thick}]
\node[box,fill=omDef!8,draw=omDef,minimum width=2.8cm] (inv) at (0,0) {investisseurs~:\\ménages, entreprises};
\node[box,fill=omThm!8,draw=omThm,minimum width=2.8cm] (mmf) at (4.6,0) {fonds\\monétaires};
\node[box,fill=omProp!8,draw=omProp,minimum width=2.8cm] (tsy) at (10.2,2.3) {Trésor\\(bons)};
\node[box,fill=omProp!8,draw=omProp,minimum width=2.8cm] (dlr) at (10.2,0.55) {dealers et banques\\(repo, CD)};
\node[box,fill=omProp!8,draw=omProp,minimum width=2.8cm] (fed) at (10.2,-1.2) {Réserve fédérale\\(reverse repo)};
\node[box,fill=gray!8,draw=gray,minimum width=2.8cm] (cp) at (10.2,-2.95) {entreprises\\(papier)};
\draw[arr] (inv) -- node[above=17pt,font=\footnotesize] {espèces contre parts} (mmf);
\draw[arr] (mmf.east) -- (tsy.west);
\draw[arr] (mmf.east) -- (dlr.west);
\draw[arr] (mmf.east) -- (fed.west);
\draw[arr,dashed] (mmf.east) -- node[below left=0pt,font=\footnotesize,align=right] {fonds prime\\uniquement} (cp.west);
\end{tikzpicture}

\omcaption{Où vont chaque jour les liquidités d'un \omterm{def:m2:money-markets:mmf}{fonds monétaire}. Un fonds gouvernemental
choisit entre les \omterm{def:m2:money-markets:bill}{bons du Trésor}, le repo avec les dealers et la \omterm{def:m2:money-markets:rrp}{facilité de reverse repo}
de la banque centrale, à leurs taux du jour~; un fonds prime peut aussi acheter du papier et des
\omterm{def:m2:money-markets:cd}{certificats de dépôt} bancaires. En fin de trimestre et à la fin de l'année, la flèche vers les dealers se réduit et
celle vers la Federal Reserve grandit (voir la \cref{fig:m2:money-markets:rrp}).}\label{fig:m2:money-markets:plumbing}
\end{omfigure}

\section{Passages de fin de période~: fin de trimestre et fin d'année}

\begin{definition}[Passage de fin de période]\label{def:m2:money-markets:turn}
Un \emph{passage de fin de période}\index{passage de fin de période} (turn) est une date à laquelle l'emprunt coûte plus cher, ou
le prêt rapporte moins que les jours voisins, parce qu'un bilan est
mesuré ce jour-là~: une fin de trimestre, et surtout la fin d'année. La
\emph{prime de passage} (turn premium) est l'écart entre le taux de cette nuit et
celui des nuits ordinaires.
\end{definition}

Les banques et les dealers sont évalués, pour leurs ratios de fonds propres, leurs notes prudentielles
et par leurs actionnaires, en partie sur la taille de leur bilan aux
dates de publication. Les repos et les dépôts sont volumineux et peu coûteux à réduire, si bien qu'ils
se réduisent à ces dates~: un dealer, qui emprunte d'ordinaire les liquidités d'un fonds au jour le jour
et les reprête, les décline pour une nuit, ou ne les accepte qu'à un prix plus élevé. Les
liquidités vont à la facilité de la banque centrale s'il en existe une, ce qui explique pourquoi la
facilité voit son recours bondir ces jours-là (voir la \cref{fig:m2:money-markets:rrp}), et
le taux qui subsiste sur le marché bondit avec lui.

\begin{omfigure}
\begin{tikzpicture}
\begin{axis}[width=11cm,height=5cm,
  ylabel={milliards de dollars},
  tick label style={font=\small},label style={font=\small},
  xmin=0,xmax=273,ymin=1900,ymax=2700,grid=major,grid style={gray!25},
  xtick={0,31,62,92,123,153,184,215,243},
  xticklabels={juil.,août,sept.,oct.,nov.,déc.,janv.,févr.,mars},
  table/col sep=comma]
\addplot[omDef,thick] table[x=t,y=usdbn]{figdata/markets-2/02-money-markets/rrp.csv};
\node[font=\footnotesize,anchor=east] at (axis cs:88,2430) {30 sept.};
\node[font=\footnotesize,anchor=east] at (axis cs:179,2560) {30 déc.~: 2\,554};
\node[font=\footnotesize,anchor=east] at (axis cs:268,2380) {31 mars};
\end{axis}
\end{tikzpicture}

\omcaption{Recours quotidien à la \omterm{def:m2:money-markets:rrp}{facilité de reverse repo} au jour le jour de la Federal Reserve,
de juillet 2022 à mars 2023. Les pics correspondent aux fins de trimestre (et,
dans une moindre mesure, certaines fins de mois)~: aux dates de déclaration, les banques et les dealers placent moins
de liquidités auprès des fonds, et la facilité absorbe le reste. Le pic de fin d'année, 245
milliards de dollars en un seul jour, est le plus élevé. Données~: Federal Reserve Bank of New York,
via la série FRED RRPONTSYD.}\label{fig:m2:money-markets:rrp}
\end{omfigure}

\begin{proposition}[Un passage de fin de période intégré dans un taux à terme]\label{prop:m2:money-markets:turn}
Un dépôt de $N$ jours au taux simple $R$ couvre $N_0$ jours ordinaires au
taux au jour le jour $r_0$, capitalisé, et un \omterm{def:m2:money-markets:turn}{passage de fin de période} de $N_T$ jours au taux
$r_T$. Si les deux manières de placer les liquidités sont équivalentes,
\[
r_T \;=\; \Bigl[\frac{1 + R\,N/B}{(1 + r_0/B)^{N_0}} - 1\Bigr]\frac{B}{N_T} .
\]
\end{proposition}
\begin{proof}
Placer au jour le jour fait passer 1 à $(1 + r_0/B)^{N_0}(1 + r_T N_T/B)$~;
le dépôt le porte à $1 + RN/B$. Égaliser les deux et résoudre en $r_T$.
\end{proof}

Un taux à terme contient donc tous les \omterm{def:m2:money-markets:turn}{passages de fin de période} qu'il couvre, dilués par sa durée~:
un passage de huit points de base sur une nuit ajoute un peu plus d'un point de base
à un taux à une semaine et un quart de point de base à un taux à un mois
(\cref{fig:m2:money-markets:oneweek}). La dilution explique aussi pourquoi un \omterm{def:m2:money-markets:turn}{passage de fin de période}
se répercute dans les taux bien avant d'arriver~: le taux à un mois coté début décembre
contient déjà la fin d'année.

\begin{omfigure}
\begin{tikzpicture}
\begin{axis}[width=11cm,height=4.8cm,
  xlabel={date de début d'un dépôt à une semaine (septembre--octobre 2026)},ylabel={taux (\%)},
  tick label style={font=\small},label style={font=\small},
  xmin=16,xmax=39,ymin=3.865,ymax=3.888,grid=major,grid style={gray!25},
  xtick={17,21,24,28,31,35,38},xticklabels={17 sept.,21,24,28,1 oct.,5,8},
  yticklabel style={/pgf/number format/fixed,/pgf/number format/precision=3},
  table/col sep=comma]
\addplot[omThm,very thick,mark=*,mark size=1.4pt] table[x=t,y=rate]{figdata/markets-2/02-money-markets/oneweek.csv};
\end{axis}
\end{tikzpicture}

\omcaption{Taux d'un dépôt à une semaine selon sa date de début, lorsque les nuits
ordinaires rapportent 3.87~\% et que la nuit de fin de trimestre du 30 septembre rapporte 3.95~\%.
Chaque dépôt qui couvre la fin de trimestre coûte 1.1 points de base de plus~; le
\omterm{def:m2:money-markets:turn}{passage de fin de période} est intégré dans le taux une semaine avant qu'il ne se produise. Taux fictifs. Données~:
le tutoriel du chapitre.}\label{fig:m2:money-markets:oneweek}
\end{omfigure}

\begin{remark}[Planchers et plafonds du passage de fin de période]\label{rem:m2:money-markets:ceiling}
Un \omterm{def:m2:central-banks-and-the-short-rate:corridor}{système de plancher} encadre le \omterm{def:m2:money-markets:turn}{passage de fin de période} des deux côtés pour ceux qui peuvent accéder aux
facilités~: un fonds peut toujours toucher le taux de reverse repo, et une banque peut
emprunter contre des Treasuries au taux de repo permanent. Aucune de ces bornes
n'est absolue. Les établissements sans accès leur échappent, et la facilité de repo
permanente ne prête que contre du collatéral et augmente le bilan de l'emprunteur
précisément la nuit où il veut afficher un bilan réduit.
\end{remark}

\section{Tutoriel~: un bon, trois taux}\label{tut:m2:money-markets:bills}

\begin{tutorial}
\textbf{Objectif.} Convertir la courbe des \omterm{def:m2:money-markets:bill}{bons du Trésor} de la \cref{fig:m2:money-markets:bills}
entre conventions, vérifier les conversions sur les exemples chiffrés du Trésor,
et calculer un \omterm{def:m2:money-markets:turn}{passage de fin de période}. \textbf{État final~:}
\cref{fig:m2:money-markets:bills} et les chiffres de
l'\cref{ex:m2:money-markets:13w}.
\begin{tutsteps}
\item \textbf{Du \omterm{def:m2:money-markets:yields}{taux d'escompte} au prix, et inversement.} Le Trésor arrondit les prix à six
décimales.
\omcode{code/firm/mmyield/firm_mmyield.py}{13}{19}{Prix d'un bon à partir de son taux d'escompte, et réciproque.}\label{lst:m2:money-markets:price}
Vérification~: 7.610~\% pour 90 jours donne 98.097500, comme dans le règlement.
\item \textbf{Le taux d'investissement}, avec sa règle des années bissextiles et sa formule quadratique
au-delà d'un semestre.
\omcode{code/firm/mmyield/firm_mmyield.py}{41}{49}{Le rendement équivalent obligataire, bons courts et longs.}\label{lst:m2:money-markets:bey}
Vérification~: un bon à 52 semaines au prix de 92.265000 donne 8.237~\%.
\item \textbf{La courbe.} \texttt{money\_market\_demo.bill\_curve()} convertit
les sept échéances~; \texttt{fig\_money\_market.py} produit
\cref{fig:m2:money-markets:bills}. La ligne à 13 semaines affiche 99.034389,
3.8572~\%, 3.9108~\%.
\item \textbf{Un \omterm{def:m2:money-markets:turn}{passage de fin de période}.}
\omcode{code/firm/mmyield/firm_mmyield.py}{57}{63}{Le taux au jour le jour sur un passage de fin de période, implicite dans un taux à terme.}\label{lst:m2:money-markets:turn}
\end{tutsteps}
\textbf{Pour aller plus loin.} Remplacer les \omterm{def:m2:money-markets:yields}{taux d'escompte} par un taux constant de 3.80~\% et
lire la courbe de \omterm{def:m2:money-markets:yields}{rendement équivalent obligataire} (\cref{exo:m2:money-markets:7})~; puis déplacer la
date d'émission au 1\up{er} mars 2027 et voir quels bons changent.
\end{tutorial}

\section{Construction~: le convertisseur de cotations du marché monétaire}\label{bld:m2:money-markets:mmyield}

\begin{build}
\bfield{Rôle} Chaque taux court que la firme miniature lit arrive dans sa
propre convention~; le constructeur de courbe (\cref{ch:m2:interest-rate-swaps}), le carnet de repo
(\cref{ch:m2:repo-and-specials}) et le desk de cash ont besoin de les ramener
à une même base. Ce composant effectue la conversion.
\bfield{Interface} \texttt{price\_from\_discount(d, days)}~;
\texttt{discount\_from\_price(p, days)}~;
\texttt{money\_market\_yield(p, days, basis)}~; \texttt{year\_days(issue)}~;
\texttt{investment\_rate(p, days, y)}~; \texttt{term\_proceeds(amount, rate,
days, basis)}~; \texttt{implied\_turn(term\_rate, term\_days, normal\_rate,
normal\_days, turn\_days)}. \bfield{Règles} Prix pour 100, arrondis à six décimales lorsqu'ils sont calculés
à partir d'un \omterm{def:m2:money-markets:yields}{taux d'escompte}~; taux exprimés en décimales~; la règle 365/366 du règlement~; la
formule des bons longs au-delà de $y/2$ jours.
\bfield{Tests d'acceptation} \texttt{code/firm/mmyield/tests/}~: les quatre exemples chiffrés
du 31 CFR 356, annexe B, à la décimale imprimée près~; la règle des années
bissextiles sur ses deux exemples~; \omterm{def:m2:money-markets:yields}{taux d'escompte} $<$ \omterm{def:m2:money-markets:yields}{taux monétaire} $<$ \omterm{def:m2:money-markets:yields}{rendement équivalent obligataire}~; la
formule longue continue à six mois~; un taux constant implique un \omterm{def:m2:money-markets:turn}{passage de fin de période} constant.
\bfield{Pour aller plus loin} Conventions des \omterm{def:m2:money-markets:cp}{billets de trésorerie} et des \omterm{def:m2:money-markets:cd}{certificats de dépôt} par devise (la livre sterling
est actual/365)~; un calendrier des \omterm{def:m2:money-markets:turn}{passages de fin de période} qui connaît chaque fin de trimestre et chaque jour férié pour
USD, EUR, GBP et JPY.
\end{build}

\begin{omsources}
\item US Treasury, \emph{Treasury Bills} (TreasuryDirect)~; 31 CFR Part 356,
annexe B, \emph{Formulas and Tables}, section VI.
\item Investment Company Institute, \emph{Money Market Fund Assets}, communiqué hebdomadaire
du 17 septembre 2026. US Securities and Exchange Commission, communiqué
de presse 2023-129, \emph{SEC Adopts Money Market Fund Reforms}, 12 juillet 2023.
\item M.\,L.\ Schapiro, témoignage sur les fonds monétaires devant le
Senate Committee on Banking, Housing and Urban Affairs, 21 juin 2012.
\item T.\,K.\ Hahn, «~Commercial paper~», Federal Reserve Bank of Richmond
\emph{Economic Quarterly} 79(2), 1993.
\item Federal Reserve Bank of New York, résultats des opérations de reverse repo
(série FRED RRPONTSYD).
\end{omsources}

\section{Exercices}

\begin{exercise}[$\star$]\label{exo:m2:money-markets:1}
Donner le prix pour 100 d'un \omterm{def:m2:money-markets:bill}{bon du Trésor} à 26~semaines (182~jours) mis en
adjudication au \omterm{def:m2:money-markets:yields}{taux d'escompte} de 3.75~\%.
\end{exercise}

\begin{exercise}[$\star$]\label{exo:m2:money-markets:2}
Un \omterm{def:m2:money-markets:bill}{bon du Trésor} à 91 jours se négocie à 99.040000. Donner son \omterm{def:m2:money-markets:yields}{taux d'escompte}, son \omterm{def:m2:money-markets:yields}{taux monétaire}
et son \omterm{def:m2:money-markets:yields}{rendement équivalent obligataire} (année de 365 jours).
\end{exercise}

\begin{exercise}[$\star$]\label{exo:m2:money-markets:3}
Une entreprise achète un \omterm{def:m2:money-markets:cd}{certificat de dépôt} de 10~millions de dollars à 90 jours, rémunéré à 4.00~\% (actual/360).
Calculer la somme qu'elle reçoit à l'échéance, puis indiquer la convention de la
\cref{def:m2:money-markets:yields} dans laquelle le taux du certificat est déjà exprimé.
\end{exercise}

\begin{exercise}[$\star\star$]\label{exo:m2:money-markets:4}
Montrer que, pour un bon à six mois au plus, le \omterm{def:m2:money-markets:yields}{rendement équivalent obligataire} dépasse toujours
le \omterm{def:m2:money-markets:yields}{taux d'escompte}. Indiquer si l'écart est plus grand pour un bon à 4~semaines ou pour un bon à
26~semaines au même \omterm{def:m2:money-markets:yields}{taux d'escompte}, et en expliquer la raison.
\end{exercise}

\begin{exercise}[$\star\star$]\label{exo:m2:money-markets:5}
Reproduire l'exemple réglementaire à 52~semaines~: un \omterm{def:m2:money-markets:yields}{taux d'escompte} de 7.65~\% pour
364 jours donne un prix de 92.265000 et un taux de placement de 8.237~\%. Expliquer
pourquoi la formule au-delà d'un semestre capitalise semestriellement.
\end{exercise}

\begin{exercise}[$\star\star$]\label{exo:m2:money-markets:6}
Un \omterm{def:m2:money-markets:mmf}{fonds monétaire} institutionnel prime de 20~milliards de dollars reçoit en un jour des demandes de rachat
pour 1.1~milliard de dollars. Préciser ce qu'impose la règle de 2023, et ce qu'un fonds gouvernemental
devrait faire.
\end{exercise}

\begin{exercise}[$\star\star\star$]\label{exo:m2:money-markets:7}
\emph{Programmation.} Avec \texttt{firm\_mmyield}, calculer le \omterm{def:m2:money-markets:yields}{rendement équivalent obligataire}
de l'ensemble des sept échéances de \omterm{def:m2:money-markets:bill}{bons du Trésor} à un \omterm{def:m2:money-markets:yields}{taux d'escompte} constant de 3.80~\% (émission
de septembre 2026). Donner les valeurs à 4~semaines et à 52~semaines, puis expliquer la
pente.
\end{exercise}

\begin{exercise}[$\star\star\star$]\label{exo:m2:money-markets:8}
\emph{Trouver l'erreur.} «~Le \omterm{def:m2:money-markets:bill}{bon du Trésor} à 13~semaines a été adjugé à 3.80~\% et celui à
52~semaines à 3.79~\%~: la courbe des bons est inversée, et le marché anticipe une
baisse de taux~». Convertir et corriger.
\end{exercise}

\section{Problème~: placer la trésorerie à cheval sur la fin d'année}

\begin{problem}[{Problème du week-end --- le passage de fin d'année implicite}]\label{pb:m2:money-markets:1}
Un trésorier d'entreprise détient 500~millions de dollars du lundi 28 décembre 2026 au
lundi 4 janvier 2027. Le jeudi 31 décembre est le dernier jour ouvré
de l'année~; le vendredi 1\up{er} janvier est un jour férié. Une banque cote un dépôt d'une semaine à
3.95~\% (actual/360). Les placements au jour le jour ordinaires sont à 3.88~\%, et la
\omterm{def:m2:money-markets:rrp}{facilité de reverse repo} de la Réserve fédérale rapporte 3.75~\% aux fonds qui peuvent
y recourir.

\medskip\noindent\textbf{Partie I --- Les cotations.}
\begin{enumerate}
\item Combien de jours dure le dépôt, et combien de jours couvre le placement
au jour le jour du 31 décembre~?
\item Quel intérêt verse le dépôt d'une semaine~?
\item Exprimer son taux sous forme de \omterm{def:m2:money-markets:yields}{rendement équivalent obligataire}.
\item Un \omterm{def:m2:money-markets:mmf}{fonds monétaire} place le même montant à la \omterm{def:m2:money-markets:rrp}{facilité de reverse repo} chaque
nuit de la semaine, en réinvestissant les intérêts. Que gagne-t-il~?
\item Lequel de ces placements comporte un risque de crédit, et envers qui~?
\end{enumerate}

\medskip\noindent\textbf{Partie II --- Le \omterm{def:m2:money-markets:turn}{passage de fin de période}.}
\begin{enumerate}[resume]
\item Combien rapporteraient les trois nuits ordinaires (28, 29 et 30 décembre) à
3.88~\%, capitalisées~?
\item Donner le taux au jour le jour du 31 décembre implicite dans la cotation
du dépôt (la \cref{prop:m2:money-markets:turn}).
\item Donner la prime de \omterm{def:m2:money-markets:turn}{passage de fin de période} en points de base.
\item Quel coût le \omterm{def:m2:money-markets:turn}{passage de fin de période} représente-t-il pour une banque qui doit emprunter
500~millions de dollars sur ce passage, par rapport aux nuits ordinaires~?
\item Pourquoi une banque accepterait-elle de le payer plutôt que de ne pas emprunter~?
\end{enumerate}

\medskip\noindent\textbf{Partie III --- Le bilan.}
\begin{enumerate}[resume]
\item Expliquer le saut observé le 30 décembre 2022 sur la \cref{fig:m2:money-markets:rrp}
à l'aide de la réponse à la question 10.
\item Un \omterm{def:m2:money-markets:mmf}{fonds monétaire} qui ne trouve aucune banque le 31 décembre place plutôt ses fonds
à la \omterm{def:m2:money-markets:rrp}{facilité de reverse repo}. Qu'abandonne-t-il sur les quatre jours, par rapport au
taux implicite du \omterm{def:m2:money-markets:turn}{passage de fin de période}~?
\item Donner le taux d'un dépôt de deux semaines à partir du 28 décembre couvrant le même
\omterm{def:m2:money-markets:turn}{passage de fin de période}, si les dix autres nuits sont ordinaires.
\item Pourquoi le \omterm{def:m2:money-markets:turn}{passage de fin de période} est-il plus petit dans le taux à deux semaines~?
\item Le \omterm{def:m2:money-markets:turn}{passage de fin de période} implicite est presque exactement égal au taux de repo permanent
de la Fed, soit 4.00~\%. Coïncidence~? Que faudrait-il attendre s'il était très supérieur~?
\end{enumerate}

\medskip\noindent\textbf{Partie IV --- Jugement.}
\begin{enumerate}[resume]
\item Pourquoi le passage de fin d'année est-il déjà présent dans les taux à un mois cotés
le 1\up{er} décembre~?
\item Citer deux facteurs qui feraient disparaître les \omterm{def:m2:money-markets:turn}{passages de fin de période}.
\item La trésorière pourrait plutôt acheter un \omterm{def:m2:money-markets:bill}{bon du Trésor} échéant le 5 janvier.
Que gagne-t-elle et que risque-t-elle~?
\item Énoncer le \emph{résultat nommé}~: le passage de fin d'année implicite et sa prime.
\item En une phrase~: pourquoi une nuit est-elle différente des autres~?
\end{enumerate}
\end{problem}

\section{Questions d'entretien}

\begin{interviewq}[$\star$ \iqroles{trader, banque}]\label{iq:m2:money-markets:1}
Pourquoi le \omterm{def:m2:money-markets:yields}{taux d'escompte} d'un bon est-il inférieur à son rendement~? Lequel utiliseriez-vous
pour le comparer à une note du Trésor à deux ans~?
\end{interviewq}

\begin{interviewq}[$\star$ \iqroles{trader, chercheur}]\label{iq:m2:money-markets:2}
Qu'est-ce qu'un passage de fin de trimestre, et pourquoi existe-t-il~?
\end{interviewq}

\begin{interviewq}[$\star\star$ \iqroles{trader, chercheur, banque}]\label{iq:m2:money-markets:3}
Pourquoi la perte subie par un \omterm{def:m2:money-markets:mmf}{fonds monétaire} en septembre 2008 a-t-elle
interrompu le marché des \omterm{def:m2:money-markets:cp}{billets de trésorerie}~?
\end{interviewq}

\begin{interviewq}[$\star\star$ \iqroles{trader}]\label{iq:m2:money-markets:4}
Un \omterm{def:m2:money-markets:mmf}{fonds monétaire} d'État peut prêter en repo à 3.85~\% ou à la Fed à 3.75~\%.
Pourquoi pourrait-il malgré tout choisir la Fed~?
\end{interviewq}

\begin{interviewq}[$\star\star$ \iqroles{développeur, trader}]\label{iq:m2:money-markets:5}
Connaissant un taux à une semaine couvrant la fin d'année et le taux des nuits
ordinaires, calculer le taux de fin d'année implicite. Quelles hypothèses avez-vous retenues~?
\end{interviewq}

\begin{interviewq}[$\star\star\star$ \iqroles{chercheur, banque}]\label{iq:m2:money-markets:6}
Comment construiriez-vous une courbe monétaire en dollars pour la première année à partir
des \omterm{def:m2:money-markets:bill}{bons du Trésor}, du repo et des swaps de taux au jour le jour (OIS)~? Que feriez-vous des \omterm{def:m2:money-markets:turn}{passages de fin de période}~?
\end{interviewq}
